% 大学物理核心公式与考点汇总 - XeLaTeX无错优化版（紧凑布局）
% 适配MiKTeX/TeXLive、Windows、中文路径
\documentclass[12pt,a4paper,oneside]{article}

% ======================== 核心宏包 ========================
\usepackage{geometry}
\usepackage{microtype}  % 优化字符间距，减少Underfull/Overfull警告
\usepackage{needspace}

% 中文与字体（修复xeCJK重定义警告：使用ctex自带配置，避免重复定义）
\usepackage{fontspec}
\setmainfont{Times New Roman}
\setsansfont{Microsoft YaHei}
\usepackage[UTF8, fontset=windows]{ctex}  % cspell:ignore ctex

% 数学与格式
\usepackage{amsmath,amssymb,amsfonts}
\usepackage{upgreek}
\usepackage{xcolor}
\usepackage{colortbl}
\usepackage{booktabs}
\usepackage{multirow}
\usepackage{makecell}
\usepackage{siunitx}
\usepackage{titlesec}
\usepackage{titletoc}
\usepackage{tcolorbox}
\usepackage{fancyhdr}
\usepackage{enumitem}

% 超链接（始终放在最后）
\usepackage[
    bookmarks=true,unicode=true,colorlinks=true,
    linkcolor=mainblue,citecolor=mainblue,urlcolor=mainblue,
    pdfborder={0 0 0}
]{hyperref}
\usepackage{hypcap}

% ======================== 基础设置 ========================
\geometry{a4paper,left=2.5cm,right=2.5cm,top=2.5cm,bottom=2.5cm,headheight=15pt}

% 颜色定义（添加% cspell:ignore mainblue 告诉拼写检查忽略该词）
\definecolor{mainblue}{RGB}{41, 90, 141}  % cspell:ignore mainblue
\definecolor{secondcyan}{RGB}{26, 117, 117}
\definecolor{thirdteal}{RGB}{27, 120, 101}
\definecolor{highlight}{RGB}{255, 240, 200}
\definecolor{tablehead}{RGB}{230, 240, 250}
\definecolor{textgray}{RGB}{60, 60, 60}

% 页眉页脚
\pagestyle{fancy}
\fancyhf{}
\fancyhead[L]{\color{mainblue}\bfseries\small 大学物理核心公式与考点汇总}
\fancyhead[R]{\color{mainblue}\thepage}
\renewcommand{\headrulewidth}{0.5pt}
\renewcommand{\headrule}{\color{mainblue}\hrule width\headwidth height\headrulewidth}
\fancypagestyle{plain}{
    \fancyhf{}
    \fancyhead[L]{\color{mainblue}\bfseries\small 大学物理核心公式与考点汇总}
    \fancyhead[R]{\color{mainblue}\thepage}
    \renewcommand{\headrulewidth}{0.5pt}
}

% 标题间距优化（紧凑）
\titlespacing*{\section}{0pt}{6pt plus 2pt minus 2pt}{3pt plus 2pt minus 2pt}
\titlespacing*{\subsection}{0pt}{4pt plus 2pt minus 2pt}{2pt plus 2pt minus 2pt}
\titlespacing*{\subsubsection}{0pt}{3pt plus 1pt minus 1pt}{1pt plus 1pt minus 1pt}

% 列表配置（紧凑）
\setlist{
    itemsep=0.1ex,
    parsep=0.2ex,
    leftmargin=1.5em,
    labelsep=0.4em,
    topsep=0.2ex,
    partopsep=0.1ex,
    font=\small
}
\setlist[itemize,2]{leftmargin=1.8em,label=$\cdot$,itemsep=0.05ex,parsep=0.1ex}

% 数学公式间距与换行优化
\numberwithin{equation}{section}
\allowdisplaybreaks
\setlength{\abovedisplayskip}{3pt}
\setlength{\belowdisplayskip}{3pt}
\setlength{\abovedisplayshortskip}{2pt}
\setlength{\belowdisplayshortskip}{2pt}
% 移除未定义的\mathindent，替换为等效的公式左对齐设置
\makeatletter
\renewcommand{\@mathmargin}{0pt}
\makeatother

% 修复Underfull/Overfull警告：调整段落和表格间距
\setlength{\emergencystretch}{2em}  % 增大拉伸空间，减少Underfull警告
\sloppy  % 放宽排版限制，减少Underfull/Overfull警告
\setlength{\parindent}{0pt}  % 取消段落缩进
\setlength{\parskip}{0.8ex}  % 适度增大段落间距，减少行填充不足
\setlength{\tabcolsep}{4pt}  % 调整表格列间距，避免内容溢出

% ======================== 自定义命令（优化换行） ========================
\newcommand{\keyformula}[1]{%
    \needspace{2\baselineskip}%
    \begin{tcolorbox}[
        colback=highlight,colframe=white,boxrule=0pt,arc=0pt,
        left=6pt,right=6pt,top=4pt,bottom=4pt,center upper,
        boxsep=0pt
    ]
    $\displaystyle #1$
    \end{tcolorbox}%
    \vspace{0.2em}
}

% 优化符号说明命令，强制换行并限制宽度
\newcommand{\symb}[2]{\textbf{#1}：\parbox[t]{10cm}{\raggedright #2}\\[0.3ex]}

\newcommand{\corepoint}[1]{%
    \begin{tcolorbox}[
        colback=white,colframe=mainblue,boxrule=0.5pt,arc=2pt,
        left=6pt,right=6pt,top=2pt,bottom=2pt,center upper,
        boxsep=0pt
    ]
    \textbf{\color{mainblue}核心考点：#1}
    \end{tcolorbox}%
    \vspace{0.2em}
}

% ======================== 文档内容 ========================
\begin{document}
% 引入封面
% 封面定义
\newcommand{\makeCover}
{
  \begin{titlepage}
    % 临时调整封面边距
    \newgeometry{left=2cm, right=2cm, top=2cm, bottom=2cm}
    
    \centering
    \vspace*{1cm}
    
    % 牛顿墓志铭（诗意改写版）
    {\Large \itshape 他以近乎神性的心智，\\以数学为炬，\\照亮行星运转、彗星行迹与海洋潮汐。\\ 艾萨克·牛顿的墓志铭}\\[0.5cm]
    \rule{\textwidth}{1pt}\\[0.5cm]
    
    % 主标题
    {\Huge \bfseries \color{mainblue} 大学物理}\\[8pt]
    {\Huge \bfseries \color{mainblue} 核心公式与考点汇总}\\[12pt]
    {\Large \itshape —— 简洁版复习讲义 ——}\\[1cm]
    
    % 分隔线
    \rule{0.6\textwidth}{0.5pt}\\[1cm]
    
    % 内容概要与适用场景
    {\Large \textbf{内容概要：} 热学 | 光学 | 量子物理}\\[8pt]
    {\Large \textbf{适用场景：} 课程复习 · 公式速查}\\[1.5cm]
    
    % 整理者信息
    {\Large \textbf{整理者：}浮萍}\\[6pt]
    {\large \textbf{联系邮箱：}\href{mailto:3071974740@qq.com}{3071974740@qq.com}}\\[6pt]
    {\large \textbf{GitHub：}\url{https://github.com/Duckweed-yhb}}\\[6pt]
    {\Large 整理日期：\today}\\[1cm]
    
    \rule{0.6\textwidth}{0.5pt}\\[1cm]
    
    % 底部信息
    {\large \textbf{版本：} v1.0 简洁版}\\[6pt]
    {\small \itshape Copyright © \the\year\ 浮萍 · 禁止商用 · 欢迎交流}\\[0.5cm]
    
    \rule{\textwidth}{1pt}
    \vfill
    
    % 恢复原始页面布局
    \restoregeometry
  \end{titlepage}
}

% 生成封面
\makeCover
\clearpage
\thispagestyle{fancy}
\tableofcontents
\clearpage

% 引入各章节内容
% 第一章：热学
\section{热学}
\subsection{理想气体基本公式}
\begin{itemize}
    \item 理想气体物态方程：\keyformula{pV=\frac{m}{M}RT}
    \item 压强微观表达式：$p = \frac{2}{3} n \overline{\varepsilon}_{\text{kt}} = \frac{2}{3} n \left( \frac{1}{2} m_0 \overline{v^2} \right)$
    \item 物态方程另一形式：$p = nkT$
    \item 玻尔兹曼常量与气体常量关系：$k=\frac{R}{N_{\text{A}}}$
\end{itemize}

\textbf{符号说明：}\\
\symb{$p$}{压强（Pa）}
\symb{$V$}{体积（m³）}
\symb{$T$}{热力学温度（K）}
\symb{$m$}{气体质量（kg）}
\symb{$M$}{摩尔质量（kg/mol）}
\symb{$n$}{分子数密度（m⁻³）}
\symb{$R$}{气体常量（8.31 J/(mol·K)）}
\symb{$k$}{玻尔兹曼常量（$1.38\times10^{-23}$ J/K）}
\symb{$m_0$}{分子质量（kg）}
\symb{$N_{\text{A}}$}{阿伏伽德罗常量（$6.02\times10^{23}$ mol⁻¹）}

\subsection{速率分布函数}
\begin{itemize}
    \item 定义：$\frac{\mathrm{d}N}{N_0} = f(v)\mathrm{d}v$ 或 $f(v) = \frac{\mathrm{d}N}{N_0 \mathrm{d}v}$
    \item 麦克斯韦速率分布：
    \keyformula{f(v) = 4\pi \left( \frac{m_0}{2\pi kT} \right)^{\frac{3}{2}} e^{-\frac{m_0 v^2}{2kT}} v^2}
    \item 归一化条件：$\int_{0}^{\infty} f(v) \mathrm{d}v = 1$
\end{itemize}

\textbf{符号说明：}\\
\symb{$f(v)$}{速率分布函数（s/m）}
\symb{$m_0$}{分子质量（kg）}
\symb{$N_0$}{总分子数（个）}

\subsection{三种统计速率}
\begin{center}
\scriptsize
\begin{tabular}{@{}p{1.8cm}|p{5cm}|p{4.8cm}@{}}
    \toprule
    \rowcolor{tablehead} 速率类型 & 公式 & 物理意义 \\
    \midrule
    最概然速率 & $v_{\text{p}} = \sqrt{\frac{2kT}{m_0}} = \sqrt{\frac{2RT}{M}}$ & 分布函数最大值对应的速率 \\
    \midrule
    平均速率 & $\bar{v} = \sqrt{\frac{8kT}{\pi m_0}} = \sqrt{\frac{8RT}{\pi M}}$ & 速率的算术平均值 \\
    \midrule
    方均根速率 & $\sqrt{\overline{v^2}} = \sqrt{\frac{3kT}{m_0}} = \sqrt{\frac{3RT}{M}}$ & 速率平方平均的平方根 \\
    \bottomrule
\end{tabular}
\end{center}

\textbf{符号说明：}\\
\symb{$v_{\text{p}}$}{最概然速率（m/s）}
\symb{$\bar{v}$}{平均速率（m/s）}
\symb{$\overline{v^2}$}{速率平方平均（m²/s²）}

\subsection{分子动能与自由度}
\begin{itemize}
    \item 平均平动动能：\keyformula{\overline{\varepsilon}_{\text{kt}} = \frac{1}{2} m_0 \overline{v^2}=\frac{3}{2}kT}
    \item 自由度表格：
        \begin{center}
            \scriptsize
            \begin{tabular}{@{}p{2.2cm}|c|c|c@{}}
                \toprule
                \rowcolor{tablehead} 分子种类 & 平动自由度 & 转动自由度 & 总自由度$i$ \\
                \midrule
                单原子分子 & 3 & 0 & 3 \\
                \midrule
                刚性双原子分子 & 3 & 2 & 5 \\
                \midrule
                刚性多原子分子 & 3 & 3 & 6 \\
                \bottomrule
            \end{tabular}
        \end{center}
    \item 分子的平均动能：$\overline{\varepsilon}_{\text{k}} = \frac{i}{2} kT$（$i$为分子的自由度数）
    \item 理想气体内能：\keyformula{E=\frac{m}{M}\frac{i}{2}RT}
\end{itemize}

\textbf{符号说明：}\\
\symb{$i$}{总自由度（无量纲）}
\symb{$E$}{内能（J）}
\symb{$\overline{\varepsilon}_{\text{k}}$}{分子平均动能（J）}

\subsection{平均自由程与碰撞频率}
\begin{itemize}
    \item 分子的碰撞截面：$\sigma = \pi d^2$
    \item 平均碰撞频率：$\bar{z} = \sqrt{2} \sigma \bar{v} n = \sqrt{2} \pi d^2 \bar{v} n$
    \item 平均自由程：
    \keyformula{\overline{\lambda} = \frac{\bar{v}}{\bar{z}} = \frac{1}{\sqrt{2} \sigma n} = \frac{kT}{\sqrt{2}\pi d^2p}}
\end{itemize}

\textbf{符号说明：}\\
\symb{$\overline{\lambda}$}{平均自由程（m）}
\symb{$\bar{z}$}{碰撞频率（Hz）}
\symb{$d$}{分子有效直径（m）}
\symb{$\sigma$}{碰撞截面（m²）}
\symb{$n$}{分子数密度（m⁻³）}

\subsection{热力学过程}
\subsubsection{热容相关}
\begin{itemize}
    \item 摩尔定容热容：$C_{V,\text{m}} = \left( \frac{\mathrm{d}Q}{\mathrm{d}T} \right)_V = \frac{\mathrm{d}E}{\mathrm{d}T} = \frac{i}{2} R$
    \item 摩尔定压热容：$C_{p,\text{m}} = \left( \frac{\mathrm{d}Q}{\mathrm{d}T} \right)_p = \frac{\mathrm{d}E}{\mathrm{d}T} + R$
    \item 迈耶公式：\keyformula{C_{p,\text{m}} = C_{V,\text{m}} + R}
    \item 比热容比：$\gamma = \frac{C_{p,\text{m}}}{C_{V,\text{m}}}$
\end{itemize}

\textbf{符号说明：}\\
\symb{$C_{V,\text{m}}$}{摩尔定容热容（J/(mol·K)）}
\symb{$C_{p,\text{m}}$}{摩尔定压热容（J/(mol·K)）}
\symb{$\gamma$}{比热容比（无量纲）}

\subsubsection{等值过程}
\begin{table}[htbp]
    \centering
    \scriptsize
    \begin{tabular}{@{}p{1.2cm}|p{1.4cm}|p{2.6cm}|p{2.6cm}|p{2.6cm}@{}}
        \toprule
        \rowcolor{tablehead} 过程类型 & 核心条件 & 做功 $W$ & 内能变化 $\Delta E$ & 热量 $Q$ \\
        \midrule
        等体过程 & $V=$常量 & $0$ & $\frac{m}{M}C_{V,\text{m}}(T_2-T_1)$ & $Q_V = \Delta E$ \\
        \midrule
        等压过程 & $p=$常量 & $p(V_2-V_1)$ & $\frac{m}{M}C_{V,\text{m}}(T_2-T_1)$ & $Q_p = \Delta E + W$ \\
        \midrule
        等温过程 & $T=$常量 & $\frac{m}{M}RT\ln\frac{V_2}{V_1}$ & $0$ & $Q_T = W$ \\
        \midrule
        绝热过程 & $Q=0$ & $\frac{1}{\gamma-1}(p_1V_1-p_2V_2)$ & $-W$ & $0$ \\
        \bottomrule
    \end{tabular}
    \caption{理想气体等值过程对比}
    \label{tab:thermo_process}
    \vspace{-0.5em}
\end{table}

\textbf{补充公式：}\\
- 绝热过程泊松公式：\\
  $TV^{\gamma-1}=C_2$；$\frac{p^{\gamma-1}}{T^\gamma}=C_3$；$pV^\gamma=C_1$\\
- 绝热线斜率：\\
  $\left( \frac{\mathrm{d}p}{\mathrm{d}V} \right)_Q = -\gamma \frac{p}{V}$ \\
  （等温线斜率：$\left( \frac{\mathrm{d}p}{\mathrm{d}V} \right)_T = -\frac{p}{V}$）

\subsection{热机与熵}
\begin{itemize}
    \item 热机效率：\keyformula{\eta = \frac{W}{Q_1} = 1 - \frac{Q_2}{Q_1}}
    \item 制冷系数：$\varepsilon = \frac{Q_2}{W} = \frac{Q_2}{Q_1 - Q_2}$
    \item 熵：$S = k\ln \Omega$
    \item 熵变：$\Delta S = k\ln \frac{\Omega_2}{\Omega_1} > 0$，$\mathrm{d}S = \frac{\mathrm{d}Q}{T}$
    \item 熵差计算：\\
      $\Delta S = \int_{1}^{2} \mathrm{d}S \geqslant \int_{1}^{2} \frac{\mathrm{d}Q}{T}$，\\
      $\Delta S = \frac{m}{M} C_{V,\text{m}} \ln\frac{T_2}{T_1} + \frac{m}{M} R \ln\frac{V_2}{V_1}$
\end{itemize}

\textbf{符号说明：}\\
\symb{$\eta$}{热机效率（无量纲）}
\symb{$\varepsilon$}{制冷系数（无量纲）}
\symb{$S$}{熵（J/K）}
\symb{$\Omega$}{热力学概率（无量纲）}
\symb{$Q_1$}{吸热（J）}
\symb{$Q_2$}{放热（J）}

\subsection{例题参考}
\textcolor{mainblue}{%
\begin{tabular}{@{}ll@{}}
热学基础： & 3、5、6、13、18、20； \\
热力学过程： & 4、5、7、8、9、10、13、15、17、19、21； \\
教材： & 《大学物理学 (第五版)》王少杰 \\
\end{tabular}
}
\clearpage
% ======================== 光学章节 - 适配主文档格式 ========================
% 子文件说明：无任何特殊宏包依赖，主文档仅需基础宏包：
% \usepackage{ctex,booktabs,array,xcolor,amsmath,amssymb}
% 引入方式：\input{光学.tex}

\section{光学}

% ======================== 半波损失（干涉基础） ========================
\subsection{半波损失（干涉核心前置知识点）}
\textbf{【核心考点】} 半波损失的判断与干涉影响 \par\vspace{0.3em}
\begin{itemize}
    \item \textbf{核心定义}：光从光疏介质（$n_1$）入射到光密介质（$n_2$）的反射光，相位突变$\pi$（等效光程差增加$\lambda/2$），称为半波损失；折射光无此现象。
    \item \textbf{产生条件（双满足）}：
          \begin{enumerate}
              \item 仅发生在\textbf{反射光}中（折射光始终无）；
              \item 入射介质折射率 $n_1 < $ 反射介质折射率 $n_2$（光疏→光密）。
          \end{enumerate}
    \item \textbf{不产生的情况}：
          \begin{enumerate}
              \item 光密→光疏介质的反射光；
              \item 任何情况下的折射光；
              \item 垂直入射（$i=0$）且$n_1/n_2$差异极小时（工程忽略）。
          \end{enumerate}
    \item \textbf{干涉影响（关键）}：
          \begin{itemize}
              \item 仅一束光有半波损失 → 干涉加强/减弱条件反转；
              \item 两束光都有/都无 → 干涉条件不变（薄膜/劈尖/牛顿环分析核心）。
          \end{itemize}
    \item \textbf{典型场景}：薄膜干涉、劈尖干涉、牛顿环干涉。
\end{itemize}

\textbf{符号说明：} \\
$\lambda$：光在真空中的波长（m）；\quad $n$：介质折射率（无量纲）；\quad $i$：入射角（rad/°）
\vspace{0.5em}

% ======================== 光的干涉 ========================
\subsection{光的干涉}
\textbf{【核心考点】} 干涉的核心是光程差与半波损失的综合判断 \par\vspace{0.3em}

\subsubsection{杨氏双缝干涉}
\begin{itemize}
    \item \textbf{实验模型}：相干光源$S_1/S_2$间距$d$，屏距$D$（$D\gg d$），屏上$P$点距中心$O$为$x$。
    \item \textbf{光程差（核心）}：
          \[
          \delta = r_2 - r_1 \approx \frac{xd}{D}
          \]
    \item \textbf{干涉条件}：
          \begin{itemize}
              \item 明纹：$\delta = \pm k\lambda \Rightarrow x = \pm k\frac{D\lambda}{d} \quad (k=0,1,2,\dots)$（$k=0$为中央明纹）；
              \item 暗纹：$\delta = \pm (2k-1)\frac{\lambda}{2} \Rightarrow x = \pm (2k-1)\frac{D\lambda}{2d} \quad (k=1,2,3,\dots)$；
              \item 条纹间距：$\Delta x = \frac{D\lambda}{d}$（等间距、明暗相间直条纹）。
          \end{itemize}
\end{itemize}

\textbf{符号说明：} \\
$d$：双缝间距（m）；\quad $D$：缝到屏的垂直距离（m）；\quad $x$：条纹到屏中心的距离（m）；\quad $k$：条纹级次（无量纲）
\vspace{0.5em}

\subsubsection{薄膜干涉（等倾干涉）}
\begin{itemize}
    \item \textbf{模型}：薄膜折射率$n_2$，两侧介质$n_1$（$n_1 < n_2$），厚度$e$，入射角$i$，折射角$\gamma$。
    \item \textbf{光程差（含半波损失）}：
          \[
          \delta = 2e\sqrt{n_2^2 - n_1^2\sin^2 i} + \frac{\lambda}{2}
          \]
    \item \textbf{干涉条件}：
          \begin{itemize}
              \item 明纹：$2e\sqrt{n_2^2 - n_1^2\sin^2 i} = \left(k - \frac{1}{2}\right)\lambda \quad (k=1,2,3,\dots)$；
              \item 暗纹：$2e\sqrt{n_2^2 - n_1^2\sin^2 i} = k\lambda \quad (k=0,1,2,\dots)$。
          \end{itemize}
    \item \textbf{核心特点}：同一入射角对应同一级条纹（等倾），条纹为同心圆环，中心级次最高。
\end{itemize}

\textbf{符号说明：} \\
$e$：薄膜厚度（m）；\quad $n_1/n_2$：介质折射率（无量纲）；\quad $i$：入射角（rad/°）
\vspace{0.5em}

\subsubsection{薄膜干涉（等厚干涉）}
\paragraph{1. 劈尖干涉（空气劈尖）}
\begin{itemize}
    \item \textbf{模型}：两玻璃片夹成空气劈尖（$n\approx1$），夹角$\theta$（极小），厚度$e$，距棱边距离$l$。
    \item \textbf{光程差（垂直入射）}：
          \[
          \delta = 2e + \frac{\lambda}{2}
          \]（上表面反射无半波损失，下表面有）
    \item \textbf{干涉条件}：
          \begin{itemize}
              \item 明纹：$2e + \frac{\lambda}{2} = k\lambda \Rightarrow e = (2k-1)\frac{\lambda}{4} \quad (k=1,2,3,\dots)$；
              \item 暗纹：$2e + \frac{\lambda}{2} = (2k+1)\frac{\lambda}{2} \Rightarrow e = k\frac{\lambda}{2} \quad (k=0,1,2,\dots)$；
              \item 棱边特征：$e=0$时为零级暗纹（标志性特征）。
          \end{itemize}
    \item \textbf{条纹间距}：
          \[
          L = \frac{\lambda}{2\theta}
          \]（$\theta$越小，条纹越稀疏）
\end{itemize}

\paragraph{2. 牛顿环（空气薄膜）}
\begin{itemize}
    \item \textbf{模型}：平凸透镜+平面玻璃，空气膜厚度$e \approx \frac{r^2}{2R}$（$R$为透镜曲率半径，$r$为环半径）。
    \item \textbf{光程差}：$\delta = 2e + \frac{\lambda}{2}$（下表面反射有半波损失）。
    \item \textbf{干涉条件}：
          \begin{itemize}
              \item 明环：$r_k = \sqrt{\left(k - \frac{1}{2}\right)\lambda R} \quad (k=1,2,3,\dots)$；
              \item 暗环：$r_k = \sqrt{k\lambda R} \quad (k=0,1,2,\dots)$；
              \item 中心特征：$r=0$时为暗斑（半波损失导致）。
          \end{itemize}
    \item \textbf{曲率半径测量（实验核心）}：
          \[
          R = \frac{d_{k+m}^2 - d_k^2}{4m\lambda}
          \]（$d_k$为第$k$圈暗环直径）
\end{itemize}

% 优化后的等厚干涉核心参数表（增加长度、优化列宽和内容）
\begin{table}[htbp]
    \centering
    \small  % 字号调整为small，提升可读性
    \renewcommand{\arraystretch}{1.5}  % 增加行高，拉长表格
    \begin{tabular}{@{}p{2.5cm}|p{4.5cm}|p{5cm}@{}}  % 优化列宽，增加表格总长度
        \toprule
        \rowcolor[RGB]{230,240,250} 干涉类型 & 核心公式（空气介质） & 关键特征与应用 \\
        \midrule
        空气劈尖 & 1. 光程差：$\delta=2e+\lambda/2$ \\ 
                  & 2. 条纹间距：$L=\lambda/(2\theta)$ \\
                  & 3. 厚度公式：$e=k\lambda/2$（暗纹） & 1. 棱边为零级暗纹（半波损失）； \\
                  &                                   & 2. 等间距平行直条纹； \\
                  &                                   & 3. 应用：测量细丝直径、玻璃平整度 \\
        \midrule
        牛顿环 & 1. 光程差：$\delta=2e+\lambda/2$ \\
                & 2. 暗环半径：$r_k=\sqrt{k\lambda R}$ \\
                & 3. 曲率半径：$R=(d_{k+m}^2-d_k^2)/(4m\lambda)$ & 1. 中心为暗斑（半波损失）； \\
                &                                             & 2. 内疏外密同心圆环； \\
                &                                             & 3. 应用：测量透镜曲率半径、波长 \\
        \bottomrule
    \end{tabular}
    \caption{等厚干涉核心参数与应用（空气介质）}
    \label{tab:equal_thickness_interference}
\end{table}
\vspace{0.5em}

% ======================== 光的衍射 ========================
\subsection{光的衍射}
\textbf{【核心考点】} 衍射的核心是半波带法（单缝）和光栅方程（光栅） \par\vspace{0.3em}

\subsubsection{惠更斯-菲涅耳原理}
\begin{itemize}
    \item \textbf{衍射现象}：光偏离直线传播，绕过障碍物形成明暗条纹的现象；
    \item \textbf{核心原理}：波前上每一点都是子波源，衍射光为所有子波的相干叠加。
\end{itemize}
\vspace{0.5em}

\subsubsection{夫琅禾费单缝衍射}
\begin{itemize}
    \item \textbf{模型}：单缝宽度$a$，入射光波长$\lambda$，衍射角$\theta$（衍射光与法线夹角）。
    \item \textbf{半波带法（核心分析方法）}：
          \begin{itemize}
              \item 最大光程差：$\Delta = a\sin\theta$；
              \item 半波带数：$N = \frac{2a\sin\theta}{\lambda}$；
              \item 叠加规律：$N$偶数→暗纹，$N$奇数→明纹（次极大），$\theta=0$→中央明纹。
          \end{itemize}
    \item \textbf{条纹公式}：
          \begin{itemize}
              \item 暗纹：
                    \[
                    a\sin\theta = \pm k\lambda \quad (k=1,2,3,\dots)
                    \]
              \item 明纹（次极大）：$a\sin\theta = \pm (2k+1)\frac{\lambda}{2} \quad (k=1,2,3,\dots)$；
              \item 中央明纹：半角宽度$\theta_0 \approx \frac{\lambda}{a}$，线宽度$\Delta x_0 \approx 2f\frac{\lambda}{a}$（$f$为透镜焦距）；
              \item 相对光强：
                    \[
                    \frac{I}{I_0} = \left( \frac{\sin\alpha}{\alpha} \right)^2
                    \]（$\alpha = \frac{\pi a\sin\theta}{\lambda}$）。
          \end{itemize}
    \item \textbf{关键结论}：$a$越小/$\lambda$越长，衍射越明显；中央明纹占80%以上能量，宽度为其他明纹的2倍。
\end{itemize}

\textbf{符号说明：} \\
$a$：单缝宽度（m）；\quad $\theta$：衍射角（rad/°）；\quad $f$：透镜焦距（m）；\quad $I_0$：中央明纹光强（W/m²）
\vspace{0.5em}

\subsubsection{光栅衍射}
\begin{itemize}
    \item \textbf{光栅常量}：$d = a + b$（$a$缝宽，$b$刻痕宽）；
    \item \textbf{光栅方程（主极大）}：
          \[
          d\sin\varphi = \pm k\lambda \quad (k=0,1,2,\dots)
          \]（所有缝光相干加强）；
    \item \textbf{缺级条件（干涉+衍射）}：
          \[
          k = \pm \frac{d}{a}k' \quad (k'=1,2,3,\dots)
          \]（主极大落在衍射暗纹处）；
    \item \textbf{核心特点}：缝数$N$越多，主极大越亮、越窄；缺级是衍射对干涉的调制结果。
\end{itemize}

\textbf{符号说明：} \\
$d$：光栅常量（m）；\quad $\varphi$：衍射角（rad/°）；\quad $k$：主极大级次（无量纲）
\vspace{0.5em}

\subsubsection{光学仪器的分辨本领}
\begin{itemize}
    \item \textbf{限制根源}：圆孔衍射形成艾里斑，无法无限分辨；
    \item \textbf{瑞利判据（临界分辨条件）}：
          \[
          \delta\theta = 1.22 \frac{\lambda}{d}
          \]（最小分辨角，$d$为透镜/圆孔直径）；
    \item \textbf{分辨本领}：$R = \frac{1}{\delta\theta} = \frac{d}{1.22\lambda}$；
    \item \textbf{提升途径}：增大$d$（天文望远镜大口径）、减小$\lambda$（电子显微镜）。
\end{itemize}
\vspace{0.5em}

\subsubsection{X射线衍射（补充）}
\textbf{【核心考点】} X射线衍射的布拉格条件 \par\vspace{0.3em}
\begin{itemize}
    \item \textbf{核心公式（布拉格条件）}：
          \[
          2d\sin\theta = k\lambda \quad (k=1,2,3,\dots)
          \]（$d$为晶面间距）；
    \item \textbf{应用}：晶体结构分析（X射线衍射仪）。
\end{itemize}
\vspace{0.5em}

% ======================== 光的偏振 ========================
\subsection{光的偏振}
\textbf{【核心考点】} 偏振是横波的特有性质，核心定律为马吕斯定律和布儒斯特定律 \par\vspace{0.3em}

\subsubsection{偏振光与自然光}
\begin{itemize}
    \item \textbf{自然光}：垂直传播方向的所有振动方向均等，可分解为两个垂直、等幅的线偏振光；
    \item \textbf{偏振光}：
          \begin{itemize}
              \item 线偏振光：仅沿一个固定方向振动；
              \item 圆/椭圆偏振光：两个垂直振动分量有固定相位差。
          \end{itemize}
\end{itemize}
\vspace{0.5em}

\subsubsection{马吕斯定律}
\begin{itemize}
    \item \textbf{自然光→线偏振光}：入射光强$I_0$，通过起偏器后光强$I_1 = \frac{I_0}{2}$；
    \item \textbf{线偏振光通过检偏器}：
          \[
          I_2 = I_1 \cos^2\alpha
          \]（$\alpha$为入射光与检偏器偏振方向的夹角）；
    \item \textbf{关键现象}：$\alpha=90^\circ$时$I_2=0$（消光），验证光的横波性。
\end{itemize}

\textbf{符号说明：} \\
$I_0$：入射自然光强（W/m²）；\quad $\alpha$：偏振方向夹角（rad/°）
\vspace{0.5em}

\subsubsection{反射和折射时光的偏振（布儒斯特定律）}
\begin{itemize}
    \item \textbf{一般情况}：反射光为垂直入射面的部分偏振光，折射光为平行入射面的部分偏振光；
    \item \textbf{布儒斯特角（起偏振角）}：
          \begin{itemize}
              \item 条件：反射光为线偏振光（仅垂直入射面振动），且反射光⊥折射光（$i_0 + r = 90^\circ$）；
              \item 公式：
                    \[
                    \tan i_0 = \frac{n_2}{n_1}
                    \]（$n_1$为入射介质折射率，$n_2$为折射介质折射率）；
              \item 示例：空气（$n_1=1$）→玻璃（$n_2=1.5$），$i_0 \approx 56^\circ$。
          \end{itemize}
    \item \textbf{玻璃片堆}：多次反射滤除垂直振动分量，折射光近似为平行入射面的线偏振光。
\end{itemize}
\vspace{0.5em}

\subsubsection{双折射现象与波片}
\paragraph{1. 双折射基础}
\begin{itemize}
    \item \textbf{寻常光（o光）}：遵守折射定律，$n_o$为定值，振动方向垂直于主平面；
    \item \textbf{非常光（e光）}：不遵守折射定律，$n_e$随传播方向变化，振动方向平行于主平面；
    \item \textbf{光轴}：晶体中o光与e光传播速度相同的方向（无双折射现象）。
\end{itemize}

\paragraph{2. 波片（偏振态转换核心元件）}
\begin{itemize}
    \item \textbf{相位差公式}：$\Delta\varphi = \frac{2\pi}{\lambda}(n_o - n_e)d$（$d$为波片厚度，$\lambda$为真空中波长）；
    \item \textbf{典型波片（对特定波长）}：
          \begin{itemize}
              \item 1/4波片：$\Delta\varphi = (2k+1)\frac{\pi}{2}$，$\alpha=45^\circ$时线偏振光→圆偏振光；
              \item 1/2波片（半波片）：$\Delta\varphi = (2k+1)\pi$，线偏振光振动面转过$2\alpha$。
          \end{itemize}
\end{itemize}

% 波片特性汇总表
\begin{table}[htbp]
    \centering
    \scriptsize
    \renewcommand{\arraystretch}{1.2}
    \begin{tabular}{@{}p{2cm}|p{3.5cm}|p{5.5cm}@{}}
        \toprule
        \rowcolor[RGB]{230,240,250} 波片类型 & 相位差/光程差 & 核心效果 \\
        \midrule
        1/4波片 & $\Delta\varphi=(2k+1)\pi/2$；$\delta=(2k+1)\lambda/4$ & 线偏振光→圆/椭圆偏振光 \\
        \midrule
        1/2波片 & $\Delta\varphi=(2k+1)\pi$；$\delta=(2k+1)\lambda/2$ & 线偏振光振动面旋转$2\alpha$ \\
        \bottomrule
    \end{tabular}
    \caption{波片核心参数与效果}
    \label{tab:waveplate}
\end{table}
\vspace{0.5em}

% ======================== 例题参考与核心考点 ========================
\subsection{例题参考与核心考点}
{
\color[RGB]{0,80,150}
\begin{tabular}{@{}p{2.2cm}p{8.8cm}@{}}
光的干涉： & 8、9、11、13、16、18、19、20（杨氏双缝、薄膜/劈尖/牛顿环）； \\
光的衍射： & 23、24、25、26、28（单缝、光栅、分辨本领）； \\
光的偏振： & 32、33、36、38、40（马吕斯定律、布儒斯特角、波片）； \\
核心考点： & 1.半波损失判断；2.干涉光程差计算；3.单缝衍射半波带法；4.光栅缺级分析；5.偏振定律应用； \\
参考教材： & 《大学物理学 (第五版)》王少杰 \\
\end{tabular}
}

% 章节末尾视觉分隔
\vspace{0.5em}
\noindent{\color[RGB]{0,80,150}\rule{\linewidth}{0.3pt}}
\vspace{0.5em}
\clearpage
% ======================== 量子物理章节 - 适配主文档格式 ========================
% 子文件说明：无任何特殊宏包依赖，主文档仅需基础宏包：
% \usepackage{ctex,booktabs,array,xcolor,amsmath,amssymb}
% 引入方式：\input{量子物理.tex}

\section{量子物理}

% ======================== 早期量子论 ========================
\subsection{早期量子论}

\subsubsection{黑体辐射和普朗克量子假设}
\textbf{【核心考点】} 黑体辐射核心定律 \par\vspace{0.3em}
\begin{itemize}
    \item \textbf{黑体模型}：能完全吸收各种波长的电磁辐射而无反射的理想物体。
    \item \textbf{单色辐出度}：单位时间从单位表面积辐射的、波长在$\lambda\sim\lambda+d\lambda$的能量与$d\lambda$之比：
          \[
          m(\lambda, T) = \frac{\mathrm{d}M}{\mathrm{d}\lambda}
          \]
    \item \textbf{辐出度}：所有波长辐射的总能量（单色辐出度的积分）：
          \[
          M(T) = \int_{0}^{\infty} m(\lambda, T) \, \mathrm{d}\lambda
          \]
    \item \textbf{斯特藩-玻尔兹曼定律}（黑体）：
          \[
          M_0(T) = \sigma_0 T^4
          \]
    \item \textbf{维恩位移定律}：
          \[
          \lambda_{\text{m}}T = b
          \]
    \item \textbf{经典物理的困难}：维恩公式（短波符合、长波偏离）、瑞利-金斯公式（长波符合、短波紫外灾难）。
    \item \textbf{普朗克量子假设}：能量量子化，谐振子能量$E = nh\nu$（$n=0,1,2,\dots$）。
\end{itemize}

\textbf{核心常数：}\\
$\sigma_0$：斯特藩常量，$\sigma_0=5.67\times10^{-8}$ W/(m²·K⁴) \\
$b$：维恩常量，$b=2.898\times10^{-3}$ m·K \\
$h$：普朗克常量，$h=6.63\times10^{-34}$ J·s
\vspace{0.5em}

\subsubsection{光电效应和爱因斯坦光子理论}
\textbf{【核心考点】} 光电效应的光子解释 \par\vspace{0.3em}
\begin{itemize}
    \item \textbf{光电效应实验规律}：饱和电流与光强成正比、存在截止频率、遏止电压与频率线性相关、响应时间极短。
    \item \textbf{光子能量}：光的粒子性体现，能量与频率成正比：
          \[
          \varepsilon = h\nu
          \]
    \item \textbf{爱因斯坦光电效应方程}：
          \[
          h\nu = \frac{1}{2}mv_{\text{m}}^2 + A_0
          \]
    \item \textbf{光的波粒二象性}：光子动量（波粒二象性核心）：
          \[
          p = \frac{h\nu}{c} = \frac{h}{\lambda}
          \]
\end{itemize}

\textbf{符号说明：}\\
$\nu$：光的频率（Hz） \\
$A_0$：金属逸出功（J） \\
$v_{\text{m}}$：光电子最大初速度（m/s） \\
$c$：真空中光速，$c=3\times10^8$ m/s
\vspace{0.5em}

\subsubsection{康普顿效应（康普顿散射）}
\textbf{【核心考点】} 康普顿效应验证光子动量 \par\vspace{0.3em}
\begin{itemize}
    \item \textbf{实验现象}：X射线散射后部分波长变长，波长改变量与散射角相关。
    \item \textbf{波长改变量公式}：
          \[
          \Delta\lambda = \lambda' - \lambda = \frac{h}{m_0 c}(1-\cos\theta) = \frac{2h}{m_0 c}\sin^2\frac{\theta}{2}
          \]
    \item \textbf{康普顿波长}（电子）：
          \[
          \lambda_{\text{C}} = \frac{h}{m_0 c} = 2.43\times10^{-12}\ \text{m}
          \]
\end{itemize}

\textbf{符号说明：}\\
$m_0$：电子静质量，$m_0=9.11\times10^{-31}$ kg \\
$\theta$：散射角（rad/°） \\
$\lambda'$：散射后波长（m）
\vspace{0.5em}

\subsubsection{氢原子光谱和玻尔理论}
\textbf{【核心考点】} 玻尔理论的量子化条件 \par\vspace{0.3em}
\begin{itemize}
    \item \textbf{经典原子模型缺陷}：汤姆孙模型（西瓜模型）、核式模型（无法解释光谱分立性）。
    \item \textbf{玻尔氢原子理论三大假设}：
          \begin{enumerate}
              \item 定态假设：原子只能处于能量不连续的定态，定态不辐射能量。
              \item 跃迁假设：原子从一个定态跃迁到另一定态，辐射/吸收光子能量$\Delta E = h\nu = |E_n - E_k|$。
              \item 轨道角动量量子化假设：
                    \[
                    L = mvr = n\frac{h}{2\pi} = n\hbar,\quad n=1,2,3,\cdots
                    \]
          \end{enumerate}
    \item \textbf{玻尔轨道半径与速度}：
          \[
          r_n = a_0 n^2,\quad v_n = \frac{\alpha c}{n}
          \]
    \item \textbf{氢原子能级公式}：
          \[
          E_n = -\frac{m e^4}{8\varepsilon_0^2 h^2 n^2} = -\frac{13.6}{n^2}\ \text{eV}
          \]
\end{itemize}

\textbf{符号说明：}\\
$\hbar$：约化普朗克常量，$\hbar = \frac{h}{2\pi} = 1.05\times10^{-34}$ J·s \\
$a_0$：玻尔半径，$a_0=0.529\times10^{-10}$ m \\
$\alpha$：精细结构常数，$\alpha\approx\frac{1}{137}$ \\
$e$：元电荷，$e=1.60\times10^{-19}$ C \\
$\varepsilon_0$：真空介电常量，$\varepsilon_0=8.85\times10^{-12}$ F/m
\vspace{0.5em}

\subsubsection{德布罗意假设与电子衍射实验}
\textbf{【核心考点】} 实物粒子的波粒二象性 \par\vspace{0.3em}
\begin{itemize}
    \item \textbf{德布罗意假设}：实物粒子具有波粒二象性（物质波/德布罗意波）：
          \[
          \boldsymbol{p} = \frac{h}{\lambda} \boldsymbol{e} = \hbar \boldsymbol{k}
          \]
    \item \textbf{德布罗意波长计算}：
          \begin{itemize}
              \item 非相对论近似（$v\ll c$）：
                    \[
                    \lambda = \frac{h}{\sqrt{2m_0 E_{\text{k}}}} = \frac{1.226}{\sqrt{U}}\ \text{nm} \quad (\text{电子，}U\text{为加速电压，单位V})
                    \]
              \item 相对论通用式：
                    \[
                    \lambda = \frac{h}{\sqrt{\left( \frac{E_{\text{k}}}{c} \right)^2 + 2m_0 E_{\text{k}}}}
                    \]
          \end{itemize}
    \item \textbf{电子衍射实验}（验证物质波）：布拉格衍射条件：
          \[
          2d\sin\varphi = k\lambda \quad (k=1,2,3,\dots)
          \]
\end{itemize}

\textbf{符号说明：}\\
$\boldsymbol{k}$：波矢，$\boldsymbol{k} = \frac{2\pi}{\lambda} \boldsymbol{e}$ \\
$E_{\text{k}}$：粒子动能（J） \\
$d$：晶面间距（m） \\
$\varphi$：入射波与晶面夹角（rad/°）
\vspace{0.5em}

% ======================== 量子力学基础 ========================
\subsection{量子力学基础}

\subsubsection{波函数的统计解释}
\textbf{【核心考点】} 波函数的概率诠释与标准条件 \par\vspace{0.3em}
\begin{itemize}
    \item \textbf{波函数}：描述微观粒子状态的函数，记为$\Psi(x,y,z,t)$（复函数）。
    \item \textbf{玻恩统计解释}：
          \begin{itemize}
              \item 概率密度：$w = |\Psi|^2 = \Psi^* \cdot \Psi$（单位体积内找到粒子的概率）。
              \item 体积元$\mathrm{d}V$内找到粒子的概率：$\mathrm{d}W = |\Psi|^2 \mathrm{d}V$。
              \item 归一化条件：$\iiint_{-\infty}^{+\infty} |\Psi|^2 \mathrm{d}x\mathrm{d}y\mathrm{d}z = 1$。
          \end{itemize}
    \item \textbf{波函数的标准条件}：单值性、有限性、连续性（概率的物理意义要求）。
\end{itemize}
\vspace{0.5em}

\subsubsection{不确定关系（测不准原理）}
\textbf{【核心考点】} 不确定关系的物理意义与应用 \par\vspace{0.3em}
\begin{itemize}
    \item \textbf{坐标-动量不确定关系}（海森伯）：
          \[
          \Delta x \Delta p_x \ge \frac{\hbar}{2}
          \]
    \item \textbf{时间-能量不确定关系}：
          \[
          \Delta E \cdot \Delta t \ge \frac{\hbar}{2}
          \]
    \item \textbf{物理意义}：微观粒子的位置与动量、能量与时间不能同时精确测量，是粒子波粒二象性的体现。
\end{itemize}

\textbf{符号说明：}\\
$\Delta x$：位置不确定量（m） \\
$\Delta p_x$：x方向动量不确定量（kg·m/s） \\
$\Delta E$：能量不确定量（J，能级宽度） \\
$\Delta t$：时间不确定量（s，能级平均寿命）
\vspace{0.5em}

\subsubsection{薛定谔方程}
\textbf{【核心考点】} 定态薛定谔方程的形式与应用 \par\vspace{0.3em}
\begin{itemize}
    \item \textbf{定态薛定谔方程}（能量不随时间变化）：
          \[
          -\frac{\hbar^2}{2m}\nabla^2 \psi + V(x,y,z)\psi = E\psi
          \]
    \item \textbf{一维定态薛定谔方程}：
          \[
          -\frac{\hbar^2}{2m}\frac{\mathrm{d}^2\psi}{\mathrm{d}x^2} + V(x)\psi = E\psi
          \]
\end{itemize}
\vspace{0.5em}

\subsubsection{一维定态问题}
\textbf{【核心考点】} 一维无限深势阱与隧道效应 \par\vspace{0.3em}
\begin{itemize}
    \item \textbf{一维无限深方势阱}（$0 < x < a$，$V=0$；$x\le0$或$x\ge a$，$V\to\infty$）：
          \begin{itemize}
              \item 归一化波函数：
                    \[
                    \psi_n(x) = \sqrt{\frac{2}{a}} \sin\left(\frac{n\pi x}{a}\right)
                    \]
              \item 能量本征值（量子化）：
                    \[
                    E_n = \frac{n^2 h^2}{8ma^2} = \frac{n^2 \pi^2 \hbar^2}{2ma^2}\quad (n=1,2,3,\dots)
                    \]
          \end{itemize}
    \item \textbf{宽度为$2a$的势阱补充}（考点）：
          \[
          E_n = \frac{n^2 h^2}{32ma^2} = \frac{n^2 \pi^2 \hbar^2}{8ma^2}\quad (n=1,2,3,\dots)
          \]
    \item \textbf{隧道效应}：微观粒子能穿透高于自身能量的势垒（应用：扫描隧穿显微镜STM）。
    \item \textbf{一维线性谐振子}：
          \begin{itemize}
              \item 势能：$V(x) = \frac{1}{2}m\omega^2 x^2$
              \item 能量本征值（零点能效应）：
                    \[
                    E_n = \left(n + \frac{1}{2}\right)\hbar\omega\quad (n=0,1,2,\dots)
                    \]
          \end{itemize}
\end{itemize}
\vspace{0.5em}

\subsubsection{原子中的电子与原子壳层结构}
\textbf{【核心考点】} 氢原子量子数与电子自旋 \par\vspace{0.3em}
\begin{itemize}
    \item \textbf{氢原子的四个量子数}（完整描述电子状态）：
          \begin{enumerate}
              \item 主量子数$n$：$n=1,2,3,\dots$，决定电子能量（主要）。
              \item 角量子数$l$：$l=0,1,\dots,n-1$，决定轨道角动量$L=\sqrt{l(l+1)}\hbar$。
              \item 磁量子数$m_l$：$m_l=0,\pm1,\dots,\pm l$，决定角动量z分量$L_z=m_l\hbar$。
              \item 自旋磁量子数$m_s$：$m_s=\pm\frac{1}{2}$，决定自旋角动量z分量$S_z=m_s\hbar$。
          \end{enumerate}
    \item \textbf{电子自旋角动量}：
          \[
          S = \sqrt{s(s+1)}\hbar,\quad s=\frac{1}{2},\quad S_z = \pm\frac{1}{2}\hbar
          \]
    \item \textbf{泡利不相容原理}：同一原子中，不可能有两个电子的四个量子数完全相同。
    \item \textbf{能量最低原理}：电子优先填充能量最低的能级（决定元素周期律）。
\end{itemize}

% 量子数汇总表格（优化排版，适配紧凑布局）
\begin{table}[htbp]
    \centering
    \scriptsize
    \renewcommand{\arraystretch}{1.5}
    \begin{tabular}{@{}p{2.2cm}|p{3cm}|p{5.6cm}@{}}
        \toprule
        \rowcolor[RGB]{230,240,250} 量子数 & 取值范围 & 物理意义 \\
        \midrule
        主量子数 $n$ & $n=1,2,3,\dots$ & 决定能量：$E_n = -\frac{13.6}{n^2}$ eV \\
        \midrule
        角量子数 $l$ & $l=0,1,\dots,n-1$ & 决定轨道角动量：$L=\sqrt{l(l+1)}\hbar$ \\
                      & （$l=0:\text{s},1:\text{p},2:\text{d},3:\text{f}$） & 决定轨道形状 \\
        \midrule
        磁量子数 $m_l$ & $m_l=0,\pm1,\dots,\pm l$ & 决定角动量z分量：$L_z=m_l\hbar$（空间量子化） \\
        \midrule
        自旋磁量子数 $m_s$ & $m_s=\pm\frac{1}{2}$ & 决定自旋角动量z分量：$S_z=m_s\hbar$ \\
        \bottomrule
    \end{tabular}
    \caption{氢原子电子的量子数汇总}
    \label{tab:quantum_number}
\end{table}
\vspace{0.5em}

% ======================== 例题参考与考点提示 ========================
\subsection{例题参考与考点提示}
{
\color[RGB]{0,80,150}
\begin{tabular}{@{}p{2cm}p{9cm}@{}}
早期量子论： & 2、4、7、8、12、14、17、18（黑体辐射、光电效应、康普顿效应、玻尔理论）； \\
量子力学： & 19、20、21、22、26、27（波函数、不确定关系、势阱、量子数）； \\
核心考点： & 1. 德布罗意波长计算；2. 不确定关系应用；3. 一维势阱能级求解；4. 氢原子量子数取值； \\
参考教材： & 《大学物理学 (第五版)》王少杰 \\
\end{tabular}
}

% 章节末尾分隔线（视觉优化）
\vspace{0.5em}
\noindent{\color[RGB]{0,80,150}\rule{\linewidth}{0.3pt}}
\vspace{0.5em}
\clearpage

\end{document}